% =========================================================================
% Capítulo 5: Mecánica Lagrangiana - Formulación Alternativa de la Dinámica
% =========================================================================
% Archivo: C5_00_introduccion.tex
% =========================================================================

\chapter{Mecánica Lagrangiana}

\section*{Introducción}

La mecánica Lagrangiana es una formulación alternativa y elegante de la mecánica clásica. En lugar de trabajar directamente con fuerzas y aceleraciones (enfoque Newtoniano), utilizamos cantidades escalares como energías y trabajamos con \textbf{coordenadas generalizadas} y el \textbf{Lagrangiano}. Este enfoque culmina el arco del curso: las restricciones que en los capítulos anteriores exigían calcular fuerzas de reacción desaparecen del planteamiento, y las ecuaciones de movimiento se obtienen por derivación sistemática de una única función escalar.

\subsection*{Ventajas de la Mecánica Lagrangiana}

\begin{enumerate}
    \item \textbf{Manejo de Restricciones}: Incorpora automáticamente restricciones de movimiento.
    \item \textbf{Coordenadas Generalizadas}: Usa solo grados de libertad independientes.
    \item \textbf{Simetría}: Revela simetrías ocultas que conducen a cantidades conservadas.
    \item \textbf{Extensibilidad}: Generaliza fácilmente a teoría de campos y mecánica cuántica.
\end{enumerate}

\section{Conceptos Fundamentales}

\subsection*{Coordenadas Generalizadas}

En lugar de usar $\{x, y, z\}$, describimos el sistema con $n$ coordenadas generalizadas $\{q_1, q_2, \ldots, q_n\}$, donde $n$ es el número de grados de libertad del sistema. La Figura \ref{fig:c5_intro_coordenadas_generalizadas} ilustra el ejemplo clásico del péndulo montado sobre un carro: aunque las dos masas viven en el plano (cuatro coordenadas cartesianas), las restricciones reducen el sistema a solo dos grados de libertad, $q_1 = x$ y $q_2 = \theta$.

% Figura: coordenadas generalizadas (péndulo sobre carro)
\begin{figure}[!ht]
    \centering
    \begin{tikzpicture}[>=Latex, scale=1.5]

        % --- CONFIGURACIÓN DE PARÁMETROS ---
        \def\xC{2.6}      % Posición del carro [u]
        \def\Wc{1.2}      % Ancho del carro [u]
        \def\Hc{0.6}      % Alto del carro [u]
        \def\Lp{1.5}      % Largo del péndulo [u]
        \def\thp{30}      % Ángulo del péndulo [deg]

        % --- SOPORTES Y ELEMENTOS FIJOS ---
        % Suelo
        \draw[thick, gris1] (0,0) -- (5.4,0);
        \fill[pattern=north east lines, pattern color=gris1!60] (0,-0.15) rectangle (5.4,0);
        % Pared izquierda
        \draw[thick, gris1] (0,0) -- (0,1.2);
        \fill[pattern=north east lines, pattern color=gris1!60] (-0.15,0) rectangle (0,1.2);

        % --- GEOMETRÍA PRINCIPAL ---
        % Resorte pared-carro
        \draw[thick, decorate, decoration={coil, segment length=4pt, amplitude=3pt, pre length=4pt, post length=4pt}]
            (0,0.55) -- (\xC-\Wc/2,0.55) node[midway, above, yshift=3pt] {$k$};
        % Carro
        \draw[fill=gris1!15, thick] (\xC-\Wc/2,0.25) rectangle (\xC+\Wc/2,0.25+\Hc);
        \node at (\xC,0.25+\Hc/2) {$M$};
        \draw[fill=gris1!40] (\xC-\Wc/4,0.15) circle (0.12);
        \draw[fill=gris1!40] (\xC+\Wc/4,0.15) circle (0.12);
        % Péndulo colgando del carro
        \coordinate (Piv) at (\xC,0.25);
        \coordinate (Mp) at ($(Piv) + (270+\thp:\Lp)$);
        \draw[thick] (Piv) -- (Mp) node[midway, right, xshift=2pt] {$L$};
        \shadedraw[ball color=gris1!30, draw=gris1] (Mp) circle (4pt);
        \node[right, xshift=3pt] at (Mp) {$m$};

        % --- VECTORES Y FUERZAS ---
        \draw[->, very thick, myred] (Mp) -- ++(0,-0.8) node[below] {$m\vec{g}$};

        % --- ANOTACIONES Y ETIQUETAS ---
        % Coordenada generalizada x
        \draw[mydashed, gris1!60] (\xC,0.25+\Hc) -- (\xC,1.5);
        \draw[<->, thick, azul] (0,1.35) -- (\xC,1.35) node[midway, above] {$x$};
        % Coordenada generalizada theta
        \draw[mydashed, gris1!60] (Piv) -- ++(0,-1.1);
        \draw[->, thick, azul] ($(Piv)+(270:0.8)$) arc (270:270+\thp:0.8);
        \node[azul] at ($(Piv)+(270+\thp/2:1.0)$) {$\theta$};
        % Grados de libertad
        \node[gris1] at (4.3,1.6) {2 grados de libertad: $q_1 = x$, $q_2 = \theta$};

    \end{tikzpicture}
    \caption{Coordenadas generalizadas del péndulo sobre un carro con resorte}
    \label{fig:c5_intro_coordenadas_generalizadas}
\end{figure}

Ejemplos:
\begin{itemize}
    \item Péndulo simple: 1 coordenada ($\theta$)
    \item Péndulo doble: 2 coordenadas ($\theta_1, \theta_2$)
    \item Cuerpo rígido en 3D: 6 coordenadas (3 de posición + 3 de rotación, o variantes)
\end{itemize}

\subsection*{El Lagrangiano}

$$L = T - V$$

donde $T$ es la energía cinética y $V$ es la energía potencial. La función $L$ es el \textbf{Lagrangiano} del sistema.

\subsection*{Ecuación de Euler-Lagrange}

Para cada coordenada generalizada $q_j$:

$$\frac{d}{dt}\left(\frac{\partial L}{\partial \dot{q}_j}\right) - \frac{\partial L}{\partial q_j} = Q_j$$

donde $Q_j$ son fuerzas generalizadas (asociadas a fuerzas no-conservativas o que no derivan de un potencial).

\section{Cantidades Conservadas}

\subsection*{Momento Canónico}

$$p_j = \frac{\partial L}{\partial \dot{q}_j}$$

Si $L$ no depende de $q_j$ (una \textbf{coordenada cíclica}), entonces $p_j$ es una cantidad conservada.

\subsection*{Función Energía (Hamiltoniano)}

$$H = \sum_j p_j \dot{q}_j - L$$

Si el Lagrangiano no depende explícitamente del tiempo ($\partial L/\partial t = 0$), entonces $H$ es constante (energía total conservada).

\section{Vínculos y Restricciones}

\subsection*{Vínculos Holonómicos}

Restricciones del tipo $f(q_1, q_2, \ldots, q_n, t) = 0$ que pueden expresarse como relaciones entre coordenadas.

Ejemplo: Péndulo con cuerda de longitud fija: $x^2 + y^2 = L^2$

\subsection*{Vínculos No-Holonómicos}

Restricciones que no pueden integrarse para obtener relaciones entre coordenadas.

Ejemplo: Rueda rodando sin deslizar.

\section{Ventajas Prácticas}

\begin{itemize}
    \item \textbf{Menos Ecuaciones}: Solo $n$ ecuaciones para $n$ grados de libertad, en lugar de $3N$ para $N$ partículas.
    \item \textbf{Ecuaciones más Simétricas}: La forma de la ecuación es la misma para cualquier elección de coordenadas.
    \item \textbf{Simetrías Explícitas}: Las simetrías conducen directamente a leyes de conservación.
\end{itemize}

\section{Ejemplos Clásicos}

\subsection*{Péndulo Simple}

$L = \frac{1}{2}m\ell^2 \dot{\theta}^2 + mg\ell \cos\theta$

Ecuación: $\ddot{\theta} + \frac{g}{\ell}\sin\theta = 0$

\subsection*{Máquina de Atwood}

Dos masas conectadas por una cuerda sobre una polea.

\subsection*{Péndulo Doble}

Sistema caótico con dinámica compleja.

\section*{Objetivos del Capítulo}

Al completar este capítulo, deberás:
\begin{itemize}
    \item Formular problemas dinámicos usando el Lagrangiano.
    \item Identificar coordenadas generalizadas y grados de libertad.
    \item Aplicar la ecuación de Euler-Lagrange.
    \item Usar simetrías para encontrar cantidades conservadas.
    \item Resolver sistemas con restricciones de movimiento.
    \item Conectar resultados Lagrangianos con análisis Newtoniano.
\end{itemize}

\newpage
